% Contenido del póster: Onorato, Pampurini y Quaranta (2024), Research in International Business and Finance. DOI 10.1016/j.ribaf.2023.102185
% Elaborado con apoyo de IA a partir del resumen oficial del artículo (OpenAlex), 2026-11-12.
\def\DOI{10.1016/j.ribaf.2023.102185}
\def\TituloEs{Eficiencia de la actividad crediticia: una comparación entre las empresas fintech y el sector bancario}
\def\TituloOriginal{Lending activity efficiency. A comparison between fintech firms and the banking sector}
\def\Autores{Onorato, Pampurini y Quaranta}
\def\Anio{2024}
\def\Revista{Research in International Business and Finance}

\def\Problema{El fenómeno fintech está cambiando el sistema financiero y la competencia en sectores como el crédito, donde tradicionalmente operaban los bancos. Falta saber si las fintech son más o menos eficientes que ellos.}
\def\Objetivo{Comparar el nivel de eficiencia de las empresas fintech que otorgan crédito con el de los bancos.}
\def\Metodo{Análisis envolvente de datos estocástico (SDEA) con datos de 2021 de todas las empresas fintech italianas dedicadas al crédito y de todos los bancos italianos.}
\def\Resultados{\begin{itemize}
  \item Los bancos muestran niveles de eficiencia más altos que las empresas fintech.
\end{itemize}}
\def\Conclusion{En el negocio del crédito, las fintech italianas todavía no alcanzan la eficiencia de los bancos, un resultado que interesa tanto a las empresas como a los reguladores.}
\def\Relevancia{En el Perú, la mayoría de las sociedades administradoras de financiamiento participativo cerraron con pérdidas en casi todos los años, y solo la plataforma líder es rentable. Sus ingresos equivalen a alrededor de \cifra{costoMediana}~\% del monto que financia. Aplicar un análisis de eficiencia como este exigiría más datos de costos de los que publica la SMV.}
\def\PalabrasClave{eficiencia, fintech, bancos, SDEA, crédito}
% --- Ficha del estudio (herramientas/posteres/ficha_articulos.py) ---
\def\Ambito{Italia, 2021 (fintech de crédito y bancos)}
\def\Enfoque{Análisis envolvente de datos estocástico (SDEA)}
\def\FiguraPeru{4}
% --- fin de la ficha ---
